The experimental design is structured to systematically investigate the properties of continuous-time signals through mathematical modeling and computational simulation. By leveraging MATLAB’s Symbolic Math Toolbox, this methodology demonstrates how theoretical signal operations can be implemented and validated within a software environment.

\subsection{Signal Construction and Transformation Methodology}
The core of this experiment involves the symbolic definition of a baseline signal, $x(t)$, which serves as the primary subject for various time-domain operations.  The signal is constructed using Heaviside step functions to isolate its behavior within the specific interval of $t = 0$ to $t = 8$ seconds. The exact symbolic math expression designed for this experiment is:

\begin{equation}
x(t) = (1 - 0.25t) \cdot (\text{heaviside}(t) - \text{heaviside}(t - 8))
\end{equation}

The experimental procedure for signal transformation required the modification of the independent variable $t$ to achieve specific outcomes:
\begin{itemize}
    \item \textbf{Time Shifting}: Replacing $t$ with $(t - 2)$ for a delay and $(t + 4)$ for an advance.
    \item \textbf{Time Reversal}: Substituting $t$ with $(-t)$.
    \item \textbf{Combined Operations}: Implementing $x(3-t)$ to achieve both flipping and shifting.
    \item \textbf{Time Scaling}: Multiplying $t$ by a constant factor of 2 ($x(2t)$) to achieve compression.
\end{itemize}

These transformations were consolidated into a single figure using a $3 \times 2$ subplot configuration to allow for simultaneous comparative analysis of the six resulting waveforms.

\subsection{Impulse Response Modeling for Filter Systems}
The methodology also addressed the representation of system characteristics through the definition of impulse response functions ($h(t)$). Four distinct filter types were modeled using combinations of the Dirac delta function ($\delta(t)$) and the Heaviside step function ($u(t)$) to simulate ideal signal processing behavior. The specific equations implemented for these models are summarized in Table 1.

\begin{table}[h]
\centering
\caption{Mathematical Expressions for Filter Impulse Responses}
\begin{tabular}{|l|l|l|}
\hline
\textbf{Filter Type} & \textbf{Symbolic Equation} & \textbf{MATLAB Implementation} \\ \hline
Lowpass & $h_1(t) = e^{-10t}u(t)$ & \texttt{exp(-10*t) * heaviside(t)} \\ \hline
Highpass & $h_2(t) = \delta(t) - e^{-10t}u(t)$ & \texttt{dirac(t) - exp(-10*t) * heaviside(t)} \\ \hline
Bandpass & $h_3(t) = e^{-10t}\cos(500t)u(t)$ & \texttt{exp(-10*t)*cos(500*t)*heaviside(t)} \\ \hline
Bandstop & $h_4(t) = \delta(t) - e^{-10t}\cos(500t)u(t)$ & \texttt{dirac(t) - exp(-10*t)*cos(500*t)*heaviside(t)} \\ \hline
\end{tabular}
\end{table}

\subsection{Numerical Computation and Statistical Correlation}
To determine the auto-correlation of the signal $x(t)$, the methodology required a transition from symbolic expressions to numerical data arrays. [cite: 193, 469, 694] This was achieved by defining a discrete time vector ranging from -10 to 10 seconds with a high-resolution step size of $dt = 0.01$. [cite: 194, 197, 470, 471, 696, 697] The symbolic expression $x(t)$ was evaluated over this vector to create the numerical array \texttt{x\_num}. [cite: 202, 472, 698] The auto-correlation was computed using the \texttt{xcorr} function, and the resulting sequence was normalized by multiplying by the time step $dt$. [cite: 218, 233, 473, 474, 699, 700]

\subsection{Experimental Pseudocode}
The following logic describes the programmatic implementation of the experiment: [cite: 32, 160, 193, 300, 432, 469, 542, 667, 694]

\begin{verbatim}
1. START
2. DEFINE symbolic variable 't'
3. CONSTRUCT signal x(t) using Heaviside functions for interval [0, 8]
4. FOR EACH transformation (shift, flip, scale):
     APPLY transformation to t
     PLOT result in 3x2 subplot grid (Figure 1)
5. CONSTRUCT signal y(t) = exp(-t)u(t) + u(-1-t)
6. PLOT y(t) for t in range [-6, 6] (Figure 2)
7. DEFINE h1(t) through h4(t) for LPF, HPF, BPF, and BSF
8. OVERLAY h1(t) and h2(t) (excluding delta) on same axis (Figure 3)
9. CONVERT symbolic x(t) to numeric array x_num with dt = 0.01
10. COMPUTE Rxx = xcorr(x_num, x_num) * dt
11. PLOT Rxx vs tau (Figure 4)
12. END
\end{verbatim}